\chapter{Your project folder}\label{ch:projectfolder}

Everything ProPrep makes for a project goes into one folder: downloaded
structures, cleaned and protonated versions, simulation input files, results
and the record of your answers. This chapter tells you what the common files
are, and, most importantly, which file each tool picks up as its input. Knowing
that saves you from building a simulation from the wrong version of your
protein.

\section{Where the project folder is}\label{sec:projectfolder-where}

The project folder is the folder named at \prompt{Project name}
(Section~\ref{sec:starting-questions}), created inside the folder you started
ProPrep from. ProPrep works inside it until you exit.

ProPrep uses the folder you started from as the project folder in these cases:
\begin{itemize}
  \item you pressed \key{Enter} at \prompt{Project name (.)};
  \item you started with \typed{--pdbid}, \typed{--pdbfile} or
  \typed{--analysis}, which skip the project question;
  \item you resumed a recorded session.
\end{itemize}
With \typed{--output-dir}, ProPrep uses (and creates) the folder you name
there (Chapter~\ref{ch:cmdline}).

\begin{tip}
To find the folder, look at the line ProPrep prints after the project name, for
example \texttt{✓ Project directory: /Users/.../ProPrep-work/myproject}. In the
browser version, \textsf{Download project} gives you a copy of it as a
\file{.zip} file.
\end{tip}

\section{A real project folder}\label{sec:projectfolder-example}

Here is what was in a project for HIV protease (\pdb{1HSG}) after loading it,
removing waters and the inhibitor (\key{4a}), assigning protonation states
(\key{4e}) and orienting it once (\key{4f}):

\begin{tabularx}{\linewidth}{@{}p{0.47\linewidth}L@{}}
\toprule
\textbf{File or folder} & \textbf{What it is} \\
\midrule
\file{1HSG.pdb} & The structure as downloaded from the Protein Data Bank. \\
\file{filtered_structure.pdb} & Protein only, made by the PDB Filter. \\
\file{filter_statistics.json} & What the filter kept and removed. \\
\file{propka_results/} & Raw results of the pKa prediction. \\
\file{ph_profile.csv} & Predicted charge over a range of pH values. \\
\file{protonation_analysis.csv} & Table of every titratable residue. \\
\file{structure_with_md_names.pdb} & Filtered protein with Amber residue names
(such as HID, HIP, ASH) for the chosen pH. \\
\file{structure_with_md_names_oriented.pdb} & The same, rotated to line up with
the x, y and z axes. \\
\file{..._oriented.orientation.json} & The rotation that was applied. \\
\file{proprep_session_20261121_122925.json} & The session file: every answer
typed. \\
\bottomrule
\end{tabularx}

\medskip
A project that went on to build and run a simulation of ubiquitin
(\pdb{1UBQ}) also held these files and folders, which the next section
explains:
\begin{itemize}
  \item \file{1UBQ_tleap.in}, \file{1UBQ.prmtop}, \file{1UBQ.rst7} and
  \file{leap.log};
  \item \file{1UBQ_parmed.pdb} and \file{1UBQ.pdb.ter_backup};
  \item the \file{simulations/} folder.
\end{itemize}

\section{The common files}\label{sec:projectfolder-files}

In the tables below, a code such as \file{1UBQ} or \file{1HRC} stands for the
name of your own structure. Appendix~\ref{app:filetypes} explains the file formats.

\subsection*{Your records}

\begin{tabularx}{\linewidth}{@{}p{0.40\linewidth}L@{}}
\toprule
\textbf{File} & \textbf{What it is for} \\
\midrule
\file{proprep_session_<date>_<time>.json} & The \term{session file}: every
answer you typed, in order. It lets you resume, undo, repeat and turn the run
into a template. Keep it. A new one is started each time you start ProPrep. \\
\file{....json.bak-<date>_<time>} & A backup of the session file, made each
time you use \key{undo} or resume. \\
\file{..._edited.json}, \file{..._template.json} & Made by the session editor
and the template creator (Chapters~\ref{ch:sessions} and~\ref{ch:batch}). \\
\file{history.json} & Only if you chose \textsf{Save command history} under
\key{w}. A list of menu actions; not needed. \\
\file{workspace.json} & Only if you chose \textsf{Save workspace} under \key{w}.
In version 1.23.0 it holds almost nothing (Section~\ref{sec:menus-workspace}). \\
\bottomrule
\end{tabularx}

\subsection*{Structures (PDB files)}

A \term{PDB file} (\file{.pdb}) is a text file listing every atom of a
structure with its position. ProPrep keeps every version it makes as a separate
file, so you can always go back to an earlier stage.

\begin{tabularx}{\linewidth}{@{}p{0.40\linewidth}p{0.10\linewidth}L@{}}
\toprule
\textbf{File} & \textbf{Made by} & \textbf{What it is} \\
\midrule
\file{1UBQ.pdb} & \key{1a} & The downloaded structure. AlphaFold and
AlphaFill models arrive as \file{.cif} files. \\
\file{1UBQ_noH.pdb} & \key{1a} & The download with hydrogen atoms
removed (if you said yes when asked). \\
\file{1HHO_biological_assembly.pdb} & \key{1b} & The biological
assembly built from the crystal's asymmetric unit. \\
\file{1HRC_aligned.pdb} & \key{3b} & A structure moved onto another
one by the Structure Aligner. \\
\file{filtered_structure.pdb} & \key{4a} & Only the chains and components you
kept. It has no header lines. \\
\file{repaired_structure.pdb} & \key{4c} & The repaired structure, renumbered
from 1. \\
\file{1HRC_redox_sites.json} & \key{2a} & The redox sites found. \\
\file{final_transformed_structure.pdb} & \key{4d} & The structure with redox
sites renamed and prepared for the force field. \\
\file{structure_with_md_names.pdb} & \key{4e} & Residue names set for the
chosen pH. Only names change; hydrogens are added later. \\
\file{name_oriented.pdb} & \key{4f} & The input structure rotated or
moved. Each further run adds another \file{_oriented}. \\
\file{bilayer.pdb} & \key{4g} & The protein packed into a lipid membrane. \\
\bottomrule
\end{tabularx}

\begin{watchout}
Several of these names are fixed, so running a tool again overwrites the
earlier file without asking. This was seen for \file{filtered_structure.pdb}
and \file{structure_with_md_names.pdb}. If you want to keep a version, copy or
rename it before running the tool again.
\end{watchout}

\subsection*{Simulation set-up and results}

\begin{tabularx}{\linewidth}{@{}p{0.40\linewidth}p{0.10\linewidth}L@{}}
\toprule
\textbf{File} & \textbf{Made by} & \textbf{What it is} \\
\midrule
\file{1UBQ_tleap.in} & \key{5a} & The instructions ProPrep wrote for
Amber's tLEaP program, which builds the simulation system. \\
\file{1UBQ.prmtop} & \key{5a} & The \term{topology}: every atom,
bond and force-field parameter of the system. \\
\file{1UBQ.rst7} & \key{5a} & The starting coordinates of every atom,
including water and ions. \\
\file{leap.log} & \key{5a} & tLEaP's full log. Look here when a build fails. \\
\file{1UBQ_parmed.pdb} & \key{5a} & A PDB file of the finished system,
in the same atom order as the topology. \\
\file{1UBQ_parmed_replay.py} & \key{5a} & A script that repeats the
final checking step by hand. \\
\file{1UBQ.pdb.ter_backup} & \key{5a} & Your PDB as it was before
ProPrep added missing chain-break (TER) lines to it. \\
\file{1UBQ.cpin} & \key{5a} & Input for constant-pH simulations. \\
\file{simulations/}\newline\file{batch_<date>_<time>/} & \key{5b} & One folder per set of
simulations, with a folder per step (minimization, heating, production) holding
its input, output, coordinates and trajectory. \\
\file{simulations/slurm_scripts/} & \key{5b} & Scripts for running on a
computing cluster. \\
\file{..._trajectory_analysis.csv} & \key{5b} & Exported analysis
results (also as \file{.json}). \\
\file{docking/} & \key{5d} & One folder per docking run, with the poses. \\
\file{small_molecule_params_MK1/} & \key{2c} & New force-field
parameters for a ligand (similar folders for metal sites and modified amino
acids). \\
\bottomrule
\end{tabularx}

\begin{watchout}
Two topology builds can write the same \file{.prmtop} name. In one test, a
failed second build left \file{1UBQ.prmtop} empty (0 bytes) and so destroyed the
good topology from the first build. Chapter~\ref{ch:topology} explains how to
avoid this.
\end{watchout}

\section{Which file each tool picks up}\label{sec:projectfolder-pickup}

ProPrep remembers every structure file it has loaded or made during the run.
When you open a tool, it has to choose one of them as its input. Each tool
chooses by \textbf{its own fixed order of preference}, not by which file you made
last. This is the most important thing to understand about the project folder.

Every tool tells you its choice on one of its first lines. Always look for it:

\begin{screen}
Using filtered structure: filtered_structure.pdb
...
Using Protonation-Updated PDB file: structure_with_md_names.pdb
...
Using Oriented PDB file: structure_with_md_names_oriented.pdb
\end{screen}

Some tools ask you instead, with a table:

\begin{screen}
          Available Structures in Workspace
┏━━━━━━━━┳━━━━━━━━━━━━━━━━━━━━━┳━━━━━━━━━━━━━━━━━━━━━━━━━━━━━━━━━━━━━━━━━━━━━━━━━━━━
┃ Option ┃ Type                ┃ File Path
┡━━━━━━━━╇━━━━━━━━━━━━━━━━━━━━━╇━━━━━━━━━━━━━━━━━━━━━━━━━━━━━━━━━━━━━━━━━━━━━━━━━━━━
│   1    │ Oriented            │ structure_with_md_names_oriented_..._oriented.pdb
│   2    │ Protonation-Updated │ structure_with_md_names.pdb
│   3    │ Filtered            │ filtered_structure.pdb
│   4    │ RCSB PDB            │ /Users/.../hsg/1HSG.pdb
└────────┴─────────────────────┴────────────────────────────────────────────────────
Select structure to use (1): §ans«4»
\end{screen}

The table below shows what each tool picked in our tests. ``First choice''
means the tool takes the first kind of file in the list that exists.

\begin{longtable}{@{}p{0.27\linewidth}p{0.67\linewidth}@{}}
\toprule
\textbf{Tool} & \textbf{Which structure it uses} \\
\midrule
\endhead
\menupath{4a}{PDB Filter} & The only structure, if there is one. If there are
several, it asks you (table above). \\
\menupath{4b}{Amino Acid Mutator} & Asks you when there are several. \\
\menupath{4c}{Structure Fixer} & The filtered structure if there is one,
otherwise the downloaded or local PDB file. Never an AlphaFold model. \\
\menupath{4e}{Protonation} & First choice: redox-prepared, then repaired, then
filtered, then hydrogen-stripped, then the download. It \textbf{never} uses an
oriented file or its own earlier result. \\
\menupath{4f}{Structure Orientator} & First choice: \textbf{an earlier oriented
file}, then protonated, then redox-prepared, then repaired, then filtered, then
the download. \\
\menupath{4d}{Redox Site Preparer} & Used \file{structure_with_md_names.pdb}
when \key{4e} had been run first, otherwise the download. \\
\menupath{5a}{Topology Generator} & Seen using, in different projects: the
membrane system (\file{bilayer.pdb}), the protonated file, the redox-prepared
file, or the download when nothing else had been made. According to the
program's source, an oriented file comes before the protonated one. \\
\menupath{1b}{Biological Assembly} & Prefers files ProPrep made itself (for
example a \file{_noH} file or an earlier assembly) over a new download. \\
\menupath{3a}{Homology Searcher} & Asks you, with a table, when several
structures are loaded. \\
\menupath{ua}{Structure Viewer} & Asks you, with a table. \\
\bottomrule
\end{longtable}

\subsection*{What this means in practice}

These orders give sensible results when you follow the usual sequence, and
surprising ones when you do not. In our tests:

\begin{itemize}
  \item \textbf{Filter, then protonate:} works. \key{4e} used
  \file{filtered_structure.pdb}.
  \item \textbf{Filter, then fix:} works. \key{4c} used the filtered structure
  and still found the missing residues listed in the original download.
  \item \textbf{Fix, then protonate:} works, but the repaired file is
  renumbered from 1. For \pdb{2RH1}, whose residues are numbered 29--1161 in the
  original file, the protonation tables then listed residues up to 430. Use the
  new numbers when you change a residue's state.
  \item \textbf{Protonate, then orient:} works. \key{4f} used the protonated
  file, and its result has the same name with \file{_oriented} added.
  \item \textbf{Orient, then filter and protonate again:} goes wrong. After a
  second filter (one chain only) and a new protonation, \key{4f} still picked the
  \emph{old} oriented file, with two chains and the old residue names, because
  an oriented file always comes first. Meanwhile \key{4e} ignored the oriented
  coordinates. The new work and the oriented file no longer matched.
  \item \textbf{Loading a second protein:} a new download is not used if ProPrep
  has already made a file of its own. After making an assembly of \pdb{1HHO}
  and then loading \pdb{4HHB}, \key{1b} still used
  \file{1HHO_biological_assembly.pdb}. Files ProPrep made always outrank a fresh
  download.
\end{itemize}

\section{Practical rules}\label{sec:projectfolder-rules}

\begin{enumerate}
  \item \textbf{One protein per project.} Start a new project (a new project
  name) for each protein or each version of a protein. Mixing proteins in one
  folder makes tools pick the wrong file.
  \item \textbf{Follow the usual order:} load, filter, fix, protonate, then
  orient, then build the topology. Chapter~\ref{ch:bigpicture} shows the full
  sequence.
  \item \textbf{Orient last.} Orientation should be the final change to your
  structure before the topology. If you change anything after orienting
  (filter, fix, mutate or protonate again), run \key{4f} again so the oriented
  file is up to date. Run it once per change: each run starts from the previous
  oriented file and adds another \file{_oriented} to the name.
  \item \textbf{Check the \texttt{Using ...} line} every time you open a tool.
  If it names the wrong file, stop and fix that first.
  \item \textbf{To switch to a different structure,} start a new project. Within
  the same run, \key{w} then \textsf{5.~Reset workspace} makes ProPrep forget
  everything it loaded and made (your files stay on disk), so you can load
  afresh.
  \item \textbf{Keep the session file.} It is the full record of how the
  structure was prepared, and the only record ProPrep can replay.
\end{enumerate}

\section{ProPrep's own folder}\label{sec:projectfolder-home}

Besides your project folders, ProPrep keeps a few things in a folder called
\file{.proprep} in your home folder (written \file{~/.proprep}). Its name starts
with a dot, so it is hidden in normal folder views. You normally never need to
open it.

\begin{tabularx}{\linewidth}{@{}p{0.36\linewidth}L@{}}
\toprule
\textbf{Inside \file{~/.proprep}} & \textbf{What it is} \\
\midrule
\file{settings.json} & Your Preferences (Section~\ref{sec:menus-prefs}),
including a GitHub token if you stored one. \\
\file{update_check.json} & When ProPrep last checked for a new version. \\
\file{web_sessions/} & Plain-text logs of browser sessions. \\
\file{forcefield_params/}, \file{transformers/} & Your personal parameter
library (Chapter~\ref{ch:forcefields}). \\
\file{library_snapshots/} & Saved copies of that library (Chapter~\ref{ch:cmdline}). \\
\file{feedback/} & Feedback that could not be sent (Chapter~\ref{ch:help}). \\
\bottomrule
\end{tabularx}
